%% kapa: exam=1 round=36 year=2025 period=2 form=A subject=appraisal_law subject_ko=감정평가관계법규 seq=34 number=34
%% source: Q-Net 감정평가사 기출문제 (cst00309) artlSeq=5231525 file=raw/1008/5231525/2025년 제36회 감정평가사 1차시험 2교시 문제지 원본.pdf
%% conversion: kordoc 4.24.51 -> tools/gen1.py (auto); review=proofread-sonnet; flags=none
\begin{question}{34}
부동산등기법령상 신탁등기에 관한 설명으로 옳지 않은 것은?
\begin{choices}
  \item 수탁자가 여러 명인 경우 등기관은 신탁재산이 공유 또는 합유인 뜻을 기록하여야 한다.
  \item 수익자나 위탁자는 수탁자를 대위하여 신탁등기를 신청할 수 있다.
  \item 여러 명의 수탁자 중 1인이 법원의 해임으로 그 임무가 종료된 경우 다른 수탁자가 여러 명일 때에는 그 전원이 공동으로 권리변경등기를 신청하여야 한다.
  \item 수탁자의 사망으로 수탁자의 임무가 종료된 경우 신수탁자는 단독으로 신탁재산에 속하는 부동산에 관한 권리이전등기를 신청할 수 있다.
  \item 신탁등기의 말소등기는 수탁자가 단독으로 신청할 수 있다.
\end{choices}
\end{question}
\answer{1}
